% Standalone subsection; requires amsmath, amssymb, booktabs, and array.
% Matches the 100k full-policy, N=M=256 paper figure and four-seed table.
\subsection{GMM40}
\label{app:gmm40}

\paragraph{Target and oracle Q-function.}
We use the fixed two-dimensional GMM40 target from the DiKL implementation,
with 40 equally weighted Gaussian components:
\begin{equation}
p_{\mathrm{GMM40}}(x)
=\frac{1}{40}\sum_{k=1}^{40}
\mathcal{N}(x;\mu_k,\sigma_{\mathrm{GT}}^2 I_2),
\qquad \sigma_{\mathrm{GT}}\approx 1.31326.
\end{equation}
The component centers are generated once using target seed~0 and remain
identical across methods and training seeds. We restrict the action domain
to $\mathcal{X}=[-40,40]^2$ and fix $Q(x)=\log p_{\mathrm{GMM40}}(x)$.
At temperature $\tau=1$, the evaluation target is therefore
\begin{equation}
\nu(x)=\frac{p_{\mathrm{GMM40}}(x)\mathbf{1}\{x\in\mathcal{X}\}}{Z_{\mathcal{X}}},
\qquad
Z_{\mathcal{X}}=\int_{\mathcal{X}}p_{\mathrm{GMM40}}(x)\,dx
\approx 0.98267.
\end{equation}
This is global truncation of the mixture, rather than separate truncation
of equally weighted components. Networks operate in normalized coordinates
$a\in[-1,1]^2$, with $x=40a$ used for oracle queries and evaluation.
No critic is learned and no ground-truth samples are provided during training;
methods use oracle values or action gradients as required by their updates.

\paragraph{Training protocol.}
We compare iBOLT, SAC, SQL, MFPO, and DIPO using 100{,}000 actor optimizer
updates, batch size~256, and training seeds $0,1,2,3$.
For iBOLT, the batch consists of 256 independent latent/proposal groups
at the same dummy state. All methods use Adam with learning rate
$3\times10^{-4}$, retaining method-specific architectures and sampling
procedures. Equal actor-update counts do not imply equal oracle-query
counts or computational cost.

\begin{table}[t]
\centering
\caption{GMM40 settings for the reported comparison.
iBOLT scales refer to normalized action coordinates $a=x/40$.
These are GMM40-specific settings, rather than the online RL defaults.}
\label{tab:gmm40_settings}
\small
\begin{tabular}{@{}p{0.50\linewidth}p{0.44\linewidth}@{}}
\toprule
Setting & Value \\
\midrule
\multicolumn{2}{@{}l}{\textit{Common settings}} \\
Target dimension / components & 2 / 40 \\
Physical action domain & $[-40,40]^2$ \\
Oracle / temperature & $Q(x)=\log p_{\mathrm{GMM40}}(x)$ / $1$ \\
Actor optimizer updates & $100{,}000$ \\
Batch size & $256$ \\
Training seeds & $0,1,2,3$ \\
Optimizer / learning rate & Adam / $3\times10^{-4}$ \\
Evaluation samples per run & $10{,}000$ \\
Samples per distribution for MMD & $2{,}048$ \\
\midrule
\multicolumn{2}{@{}l}{\textit{iBOLT settings}} \\
Actor hidden layers / activation & $(256,256,256)$ / GELU \\
Latent prior & $z\sim\mathcal{N}(0,I_2)$, freshly sampled \\
Mixture components per group, $N$ & $256$ \\
Reference actions per group, $M$ & $256$ \\
Conditional distribution & Diagonal box-truncated Gaussian \\
Conditional center & $\tanh(\mu_\theta(z))$ \\
Actor log-standard-deviation bounds & $[-5,-3.5]$ \\
Initial log-standard deviation & $-4$ \\
Proposal-only standard-deviation floor & $0.05$ \\
Mean-head initialization variance scale & $16$ (fan-average uniform) \\
Log-scale head initial weights / bias & $0$ / $-4$ \\
Density-correction exponent & $1$ \\
Gradient clipping & None \\
\bottomrule
\end{tabular}
\end{table}

\paragraph{iBOLT proposal and initialization.}
Each update draws fresh latent variables and constructs the finite mixture
used in the actor likelihood. To broaden candidate sampling, the proposal
uses $\widetilde{\sigma}_{\theta,d}(z)=\max(\sigma_{\theta,d}(z),0.05)$.
Both proposal sampling and its density $q(a)$ use these broadened scales,
while the actor likelihood and final policy sampling use the original
actor scales. Thus, in this experiment, the proposal is distinct from the
finite mixture used in the likelihood. Reference actions are sampled by
uniformly choosing mixture components with replacement and drawing from
their truncated Gaussians. The importance weights satisfy
$w_j\propto\exp(Q(40a_j))/q(a_j)$ and are normalized within each group.
Reference actions and weights are detached during the likelihood update.
The mean head uses a fan-average uniform initializer with variance scale~16;
this scales the initialization variance, not the output by a factor of~16.

\paragraph{Baseline adaptations.}
SAC uses a $(256,256)$ ReLU actor and a tanh-squashed diagonal Gaussian,
with log-standard-deviation bounds $[-5,2]$ and initial standard deviation
$0.5$ before squashing. Its actor minimizes
$\mathbb{E}[\log\pi(x)-Q(x)]$ with the exact oracle.
SQL uses a $(256,256)$ ReLU implicit actor and amortized SVGD with
256 particles, split into 128 fixed and 128 updated particles.
Its adaptive RBF kernel is $k(a,b)=\exp(-\|a-b\|^2/h)$, where
$h=\max(\operatorname{median}_{\mathrm{lower}}\|a-b\|^2/\log 128,10^{-3})$,
computed over fixed--updated particle pairs in normalized coordinates.
The JAX port retains the upstream bounded-action score and
TensorFlow-1-compatible Adam convention.
DIPO uses 100 diffusion steps with a cosine noise schedule and 20 action
refinement steps per actor update. Action refinement uses Adam with learning
rate~$0.03$ and gradient-norm limit~$0.2$; the actor gradient-norm limit is~$2$.
Both optimizers use $\epsilon=10^{-5}$. Its action buffer has capacity
$10^6$ and starts with 5{,}000 uniform candidates; 32 new policy samples
are added every 32 actor updates.
MFPO uses separate mean-flow and divergence networks with three
256-unit hidden layers, GELU activations, and layer normalization.
Each actor update uses 16 policy candidates and 32 additional candidates;
the divergence network is also updated once. Policy generation uses two steps.

\paragraph{Policy sampling and aggregation.}
At the final checkpoint, we draw 10{,}000 samples from each learned policy.
iBOLT uses fresh Gaussian latents and conditional truncated-Gaussian noise;
the reported results are not mean-only outputs. Baselines retain their
stochastic generators, including reverse-diffusion noise for DIPO.
We draw an independent 10{,}000-sample reference from $\nu$ using seed
20260917 and share it across methods and training seeds.
The main figure uses training seed~2 for every method.
Reported numerical results are the mean and sample standard deviation
across the four training seeds, computed from per-seed metrics.

\paragraph{Mode coverage.}
We measure coverage of the 40 reference components, rather than count
local maxima of the learned density. A sample $x$ is assigned to
$k^*(x)=\arg\min_k d_k^2(x)$, where
$d_k^2(x)=\|x-\mu_k\|^2/\sigma_{\mathrm{GT}}^2$,
and contributes to the assigned component only if $d_{k^*(x)}^2(x)\leq 9$.
Let $c_k$ and $c_k^{\mathrm{GT}}$ denote the resulting counts from the
10{,}000 policy and reference samples, respectively. We report
\begin{equation}
\operatorname{Coverage}
=\sum_{k=1}^{40}\mathbf{1}\!\left\{
c_k\geq\max(10,0.1c_k^{\mathrm{GT}})\right\}.
\end{equation}
Thus, an isolated sample near a component does not suffice for coverage,
and overlapping neighborhoods do not cause a sample to be counted twice.

\paragraph{Maximum mean discrepancy.}
We use the first 2{,}048 policy samples and the first 2{,}048 reference
samples in physical coordinates, with the shared kernel
\begin{equation}
k(x,y)=\frac{1}{5}\sum_{h\in\{1,2,5,10,20\}}
\exp\!\left(-\frac{\|x-y\|^2}{2h^2}\right).
\end{equation}
We compute the unbiased U-statistic estimator of squared MMD, excluding
diagonal terms within each sample set, and report
$\sqrt{\max(\widehat{\operatorname{MMD}}_u^2,0)}$.
The square root is taken separately for each seed before aggregation.

\paragraph{Implementation.}
For this fixed-Q comparison, iBOLT, SQL, and MFPO use JAX, while SAC and
DIPO use PyTorch. The reported iBOLT runs use an RTX~4090; the SAC, MFPO,
and DIPO runs use RTX~5090 GPUs. These benchmark-specific implementations
and settings should be distinguished from the online RL implementations.
